\documentclass{article}
\usepackage{amsmath, amssymb, amsthm}
\usepackage{hyperref}
\usepackage{geometry}
\geometry{margin=1in}

\title{Erd\H{o}s Problem \#1208}
\date{Last updated: 2026-04-04}
\author{Source: erdosproblems.com}

\begin{document}
\maketitle

\noindent\textbf{Status:} OPEN \quad \textbf{No prize}

\noindent\textbf{Tags:} geometry, distances

\medskip
\noindent\textbf{Problem Statement:}

For $d\geq 2$ let $F_d(n)$ be minimal such that every set of $n$ points in $\mathbb{R}^d$ contains a set of $F_d(n)$ points with distinct distances. Estimate $F_d(n)$ for fixed $d$ as $n\to \infty$.

\bigskip
\noindent\textbf{Remarks:}

It is known that\[\frac{n^{1/3}}{(\log n)^{1/3}}\ll F_2(n) \ll \frac{n^{1/2}}{(\log n)^{1/4}}.\]The upper bound is demonstrated by the set of integer lattice points in a square grid (which have $\ll n/\sqrt{\log n}$ many distinct distances). The lower bound was proved by Charalambides \cite{Ch13}.

Thiele \cite{Th95} proved $F_d(n) \gg n^{\frac{1}{3d-2}}$ for all $d\geq 3$. This was improved to\[F_d(n) \gg_d n^{\frac{1}{3d-3}}(\log n)^{1/3-\frac{2}{3d-3}}\]by Conlon, Fox, Gasarch, Harris, Ulrich, and Zbarsky \cite{CFGHUZ15}. The example of integer lattice points shows that $F_d(n)\ll n^{1/d}$ for $d\geq 2$.

Erd\H{o}s and Guy \cite{ErGu70} consider the problem of determining the maximal size of such a subset of the $n^{1/d}\times\cdots \times n^{1/d}$ integer grid, showing that it is at least $n^{\frac{2}{3d}-o(1)}$ (and at most $n^{1/d}$). Lefmann and Thiele \cite{LeTh95} have improved the lower bound to $n^{\frac{2}{3d}}$.

The case $d=1$ is the subject of [530], and Koml\'{o}s, Sulyok, and Szemer\'{e}di \cite{KSS75} have proved that $F_1(n)\asymp n^{1/2}$.

See also [1207]. The behaviour of $F_d(n)$ for fixed $n\geq 3$ as $d\to \infty$ is the subject of [1088].

\bigskip
\noindent\textbf{References:}

[CFGHUZ15] Conlon, David and Fox, Jacob and Gasarch, William and Harris,
David G. and Ulrich, Douglas and Zbarsky, Samuel, Distinct volume subsets. SIAM J. Discrete Math. (2015), 472--480.
 
 [Ch13] Charalambides, Marcos, A note on distinct distance subsets. J. Geom. (2013), 439--442.
 
 [ErGu70] Erd\H{o}s, P. and Guy, R. K., Distinct distances between lattice points. Elem. Math. (1970), 121--123.
 
 [KSS75] Koml\'{o}s, J. and Sulyok, M. and Szemeredi, E., Linear problems in combinatorial number theory. Acta Math. Acad. Sci. Hungar. (1975), 113-121.
 
 [LeTh95] Lefmann, Hanno and Thiele, Torsten, Point sets with distinct distances. Combinatorica (1995), 379--408.
 
 [Th95] T. Thiele, Geometric selection problems and hypergraphs. PhD thesis, Freir Universitat Berlin (1995).

\bigskip
\noindent\small{Source: \url{https://www.erdosproblems.com/1208}}
\end{document}
